\section{Programm: Reihen}
\label{sec:c:reihen}

Die Zahl $\pi$ ist eine sehr wichtige Zahl in der Mathematik. Sie lässt sich aus einer unendlichen Reihe berechnen:
\begin{equation*}
  \frac{\pi}{4} = 1 - \frac{1}{3} + \frac{1}{5} - \frac{1}{7} + \frac{1}{9} - \frac{1}{11} + \ldots
\end{equation*}
Die Reihe geht also unendlich weiter. In der Praxis kann man aber nur endlich viele Terme berechnen. Je mehr Terme man berechnet, desto genauer wird das Ergebnis. Im folgenden Programm wird diese Summe mit $1000000$ Schritten berechnet (auf der Konsole wird $\pi$ ausgegeben, indem man obige Summe einfach vervierfacht):

\begin{fullwidth}
  \cppinput{c/04_reihe.c}  
\end{fullwidth}

\begin{exercise}
  Einige von Euch haben sich vielleicht schon ein paar Stellen von $\pi$ gemerkt und werden feststellen, dass das Ergebnis sechs Stellen nach dem Komma abweicht. Ihr könnt einmal das \texttt{printf()} mit folgernder Anweisung
  tauschen, die Euch die Abweichung von $\pi$ anzeigt:
\begin{code}
printf("%9d Pi: %.15f  (% .15f)\n", i, 4*pi, 4*pi - 3.14159265358979323846);
\end{code}
  Wie groß wäre die Abweichung nach 1000 Schritten?
  \begin{solution}
    Ab der 7. Programmzeile wird
\begin{code}
double n = 2*i + 1  
n = 1.0/n;
\end{code}
  berechnet. Dieses $n$ wird dann zu dem Ergebnis hinzugezäht oder abgezogen – damit ist die Abweichung für $i = 1000$:
\begin{equation*}
  \Delta\pi = 4\frac{1}{2 n + 1} = \frac{4}{2001} \approx 0,002\,.
\end{equation*}
  \end{solution}
\end{exercise}

% \paragraph{Erklärung:}

% \begin{description}
%   \item[\texttt{int a = 0;}] \hfill\\
%     Wir initialisieren die erste Fibonacci-Zahl mit 0.
%   \item[\texttt{int b = 1;}] \hfill\\
%     Wir initialisieren die zweite Fibonacci-Zahl mit 1.
%   \item[\texttt{for(int i = 1; i <= 16; i = i+1) \{ … \} }] \hfill\\
%     Wir starten eine Schleife, die 16 Mal durchläuft, um die ersten 16 Fibonacci-Zahlen zu berechnen.
%     Seht auch mal den Abschnitt \ref{sec:referenz:c:schleifen} in der Referenz, um mehr über Schleifen zu erfahren.
%   \item[\texttt{int next = a + b;}] \hfill\\
%     Wir berechnen die nächste Fibonacci-Zahl, indem wir die beiden vorhergehenden Zahlen addieren.
%   \item[\texttt{a = b;}] \hfill\\
%     Wir verschieben die Werte: Die aktuelle zweite Zahl wird zur ersten Zahl.
%   \item[\texttt{b = next;}] \hfill\\
%     Wir verschieben die Werte: Die berechnete nächste Zahl wird zur zweiten Zahl.
%   \item[\texttt{printf("Fibonacci \%d: \%d$\backslash$n", i, a);}] \hfill\\
%     Wir geben die aktuelle Fibonacci-Zahl auf dem Bildschirm aus. Das \texttt{\%d} ist ein Platzhalter für die aktuelle Zahl, die an dieser Stelle ausgegeben wird. 
% \end{description}


% \paragraph{Merke:}

% \begin{itemize}[noitemsep]
%     \item Ein C-Programm beginnt immer in der \texttt{main()}-Funktion.
%     \item \texttt{printf()} schreibt Text auf den Bildschirm.
%     \item \texttt{$\backslash$n} macht einen Zeilenumbruch.
%     \item \texttt{return 0;} sagt dem Computer: ``Alles OK.''
%     \item Jede Anweisung endet mit \texttt{;}.
% \end{itemize}
